% GENERATED FILE. Edit canonical lessons, then run npm run generate:books.
% canonical-source-hash: fnv1a64:50fe55e46f2b71ff
% canonical-lessons: MR-C206-sarvajanik, MR-C206-khajgi, MR-C206-vaiyaktik, MR-C206-samanya, MR-C206-khas

\chapter{Public, Private, Personal, Ordinary, Special}
\label{ch:mr-public-private-personal-ordinary-special}
\hlchaptermodality{206}

\begin{chapteropening}
\textbf{By the end of this chapter:} \emph{I can say public, private, personal, ordinary and special.}

Complete the last lesson of chapter 206: I can say public, private, personal, ordinary and special.
\end{chapteropening}

\section[\texorpdfstring{\mr{सार्वजनिक} (sārvajanik)}{सार्वजनिक (sārvajanik)}]{\mr{सार्वजनिक} (sārvajanik) — public}

\label{lesson:MR-C206-sarvajanik}



\begin{quote}
Before the new one: say the Marathi for cooked, then the Marathi for ripe.
\end{quote}

\subsection*{You'll want to know: \mr{सार्वजनिक}}
\textbf{\mr{सार्वजनिक}} — \emph{sārvajanik} — \textquotedblleft{}public\textquotedblright{}.

\begin{grammarlens}[title={What you've built}]
One more word for this chapter.
\end{grammarlens}

\subsection*{Your turn}
\begin{itemize}
\raggedright
  \item \emph{Say it:} \emph{sārvajanik}
  \item \emph{Say it:} \emph{sārvajanik}, once more
  \item \emph{Say it:} say \emph{sārvajanik}
  \item \emph{Recall:} say the Marathi for cooked, then the Marathi for ripe, then say \emph{sārvajanik} again
\end{itemize}

\begin{tcolorbox}[breakable,colback=teal!4,colframe=teal!35!black,title={Before you move on}]
What does \mr{सार्वजनिक} mean? (\textquotedblleft{}Public\textquotedblright{}.) Say it once more.
\end{tcolorbox}
\section[\texorpdfstring{\mr{खाजगी} (khājgī)}{खाजगी (khājgī)}]{\mr{खाजगी} (khājgī) — private}

\label{lesson:MR-C206-khajgi}



\begin{quote}
Before the new one: say the Marathi for ripe, then the Marathi for public.
\end{quote}

\subsection*{You'll want to know: \mr{खाजगी}}
\textbf{\mr{खाजगी}} — \emph{khājgī} — \textquotedblleft{}private\textquotedblright{}.

\begin{grammarlens}[title={What you've built}]
One more word for this chapter.
\end{grammarlens}

\subsection*{Your turn}
\begin{itemize}
\raggedright
  \item \emph{Say it:} \emph{khājgī}
  \item \emph{Say it:} \emph{khājgī}, once more
  \item \emph{Say it:} say \emph{khājgī}
  \item \emph{Recall:} say the Marathi for ripe, then the Marathi for public, then say \emph{khājgī} again
\end{itemize}

\begin{tcolorbox}[breakable,colback=teal!4,colframe=teal!35!black,title={Before you move on}]
What does \mr{खाजगी} mean? (\textquotedblleft{}Private\textquotedblright{}.) Say it once more.
\end{tcolorbox}
\section[\texorpdfstring{\mr{वैयक्तिक} (vaiyaktik)}{वैयक्तिक (vaiyaktik)}]{\mr{वैयक्तिक} (vaiyaktik) — personal}

\label{lesson:MR-C206-vaiyaktik}



\begin{quote}
Before the new one: say the Marathi for public, then the Marathi for private.
\end{quote}

\subsection*{You'll want to know: \mr{वैयक्तिक}}
\textbf{\mr{वैयक्तिक}} — \emph{vaiyaktik} — \textquotedblleft{}personal\textquotedblright{}.

\begin{grammarlens}[title={What you've built}]
One more word for this chapter.
\end{grammarlens}

\subsection*{Your turn}
\begin{itemize}
\raggedright
  \item \emph{Say it:} \emph{vaiyaktik}
  \item \emph{Say it:} \emph{vaiyaktik}, once more
  \item \emph{Say it:} say \emph{vaiyaktik}
  \item \emph{Recall:} say the Marathi for public, then the Marathi for private, then say \emph{vaiyaktik} again
\end{itemize}

\begin{tcolorbox}[breakable,colback=teal!4,colframe=teal!35!black,title={Before you move on}]
What does \mr{वैयक्तिक} mean? (\textquotedblleft{}Personal\textquotedblright{}.) Say it once more.
\end{tcolorbox}
\section[\texorpdfstring{\mr{सामान्य} (sāmānya)}{सामान्य (sāmānya)}]{\mr{सामान्य} (sāmānya) — ordinary, common}

\label{lesson:MR-C206-samanya}



\begin{quote}
Before the new one: say the Marathi for private, then the Marathi for personal.
\end{quote}

\subsection*{You'll want to know: \mr{सामान्य}}
\textbf{\mr{सामान्य}} — \emph{sāmānya} — \textquotedblleft{}ordinary, common\textquotedblright{}.

\begin{grammarlens}[title={What you've built}]
One more word for this chapter.
\end{grammarlens}

\subsection*{Your turn}
\begin{itemize}
\raggedright
  \item \emph{Say it:} \emph{sāmānya}
  \item \emph{Say it:} \emph{sāmānya}, once more
  \item \emph{Say it:} say \emph{sāmānya}
  \item \emph{Recall:} say the Marathi for private, then the Marathi for personal, then say \emph{sāmānya} again
\end{itemize}

\begin{tcolorbox}[breakable,colback=teal!4,colframe=teal!35!black,title={Before you move on}]
What does \mr{सामान्य} mean? (\textquotedblleft{}Ordinary, common\textquotedblright{}.) Say it once more.
\end{tcolorbox}
\section[\texorpdfstring{\mr{खास} (khās)}{खास (khās)}]{\mr{खास} (khās) — special}

\label{lesson:MR-C206-khas}



\begin{quote}
Before the new one: say the Marathi for personal, then the Marathi for ordinary.
\end{quote}

\subsection*{You'll want to know: \mr{खास}}
\textbf{\mr{खास}} — \emph{khās} — \textquotedblleft{}special\textquotedblright{}.

\begin{grammarlens}[title={What you've built}]
One more word for this chapter.
\end{grammarlens}

\subsection*{Your turn}
\begin{itemize}
\raggedright
  \item \emph{Say it:} \emph{khās}
  \item \emph{Say it:} \emph{khās}, once more
  \item \emph{Say it:} say \emph{khās}
  \item \emph{Recall:} say the Marathi for personal, then the Marathi for ordinary, then say \emph{khās} again
\end{itemize}

\begin{tcolorbox}[breakable,colback=teal!4,colframe=teal!35!black,title={Before you move on}]
What does \mr{खास} mean? (\textquotedblleft{}Special\textquotedblright{}.) Say it once more.
\end{tcolorbox}
